\section{Introduction}\label{sec:introduction}

The good minimal model problem asks whether a pseudo-effective
canonical adjoint can be made nef and generated by its sections on a
birational model. It joins two different questions: construction of a
minimal model and abundance of its nef adjoint. In the compact
K\"ahler category, the existence of the first pluricanonical section
is itself a substantial part of the problem. Positivity is analytic,
and even the passage from a cohomological identity to an identity of
holomorphic line bundles must be kept explicit.

The projective minimal model program provides the birational framework
\cite{BirationalGeometry,BCHM}. In dimension three the K\"ahler
theory has developed through nonvanishing, minimal models and
abundance \cite{DemaillyPeternellNonvanishing,HPMinimalModels,
CHPAbundance,CHPAbundanceErratum,GuenanciaPaunBogomolovGieseker}.
Work on fourfolds and on
pseudo-effective adjoints supplies further analytic and birational
tools \cite{DHPFourfoldMMP,HLLNonvanishing}. Our concern is the
precise additional nonvanishing argument needed when a
pseudo-effective fourfold MMP and abundance after nonvanishing are
available.

We make three explicit hypotheses, stated in
Section~\ref{sec:premises}: full orbifold Iitaka subadditivity in
Fujiki class \(\mathcal C\), the specified pseudo-effective klt
fourfold MMP conditional on that subadditivity, and abundance for nef
klt fourfold adjoints of nonnegative Kodaira dimension. Under these
premises we resolve positively the good-minimal-model assertion in
the following category.

\begin{theorem}[Conditional good minimal models]\label{thm:main}
Assume Assumptions~\ref{hyp:campana}, \ref{hyp:mmp}, and
\ref{hyp:effective-abundance}. Let \(X\) be a normal connected
compact K\"ahler fourfold, and let \(B\) be an effective rational
Weil divisor. Suppose that \((X,B)\) is klt, that its actual adjoint
\(D=K_X+B\) is \(\Q\)-Cartier and analytically
pseudo-effective, and that \(X\) is globally strongly
\(\Q\)-factorial in the sense of Definition~\ref{def:strong}.

Then there are a normal connected globally strongly
\(\Q\)-factorial compact K\"ahler fourfold \(Y\) and a
bimeromorphic map \(\phi:X\dashrightarrow Y\) such that:
\begin{enumerate}[label=\textup{(\roman*)}]
\item \(\phi\) extracts no prime divisor, and
      \(B_Y=\phi_*B\);
\item \((Y,B_Y)\) is klt and its actual adjoint
      \(D_Y=K_Y+B_Y\) is \(\Q\)-Cartier and analytically nef;
\item \(a(E;Y,B_Y)\ge a(E;X,B)\) for every prime divisor over
      the models, with strict inequality for each prime on \(X\)
      contracted by \(\phi\);
\item for some \(m>0\), the Cartier line bundle
      \(\OO_Y(mD_Y)\) is globally generated.
\end{enumerate}
\end{theorem}

The exponent may depend on the pair. Zero boundary, smooth fourfolds
and all numerical dimensions are included. Global strong
\(\Q\)-factoriality concerns all rank-one reflexive sheaves on
the whole space; it is not an assumption of local analytic
\(\Q\)-factoriality on every open subset. The conclusion concerns
the actual holomorphic adjoint line on the specified endpoint.

Assumption~\ref{hyp:mmp} already produces a nef minimal model with
properties~\textup{(i)}--\textup{(iii)}. The new conclusion needed
before Assumption~\ref{hyp:effective-abundance} can be applied is
the following.

\begin{theorem}[Nonvanishing on the nef endpoint]
\label{thm:nonvanishing}
Under Assumptions~\ref{hyp:campana}, \ref{hyp:mmp}, and
\ref{hyp:effective-abundance}, let \((Y,\Delta)\) be an ordinary
klt pair on a normal connected globally strongly
\(\Q\)-factorial compact K\"ahler fourfold. If \(\Delta\)
is effective and rational and the actual \(\Q\)-Cartier adjoint
\(J=K_Y+\Delta\) is analytically nef, then
\(\kappa(Y,J)\ge0\).
\end{theorem}

\subsection{Structure of the proof}

The projective case uses conditional projective log abundance. We
include its complete proof in
Appendices~\ref{proj:sec:introduction}--\ref{proj:sec:completion};
Lemma~\ref{red:log-iitaka} derives its precise logarithmic-Iitaka
premise from Assumption~\ref{hyp:campana}. Thus the projective
branch introduces no further conditional premise. For a
nonprojective endpoint, rational quotients, the Albanese map and
lower-dimensional nonvanishing reduce the problem to canonical
nonvanishing on a smooth non-uniruled fourfold, principally with
irregularity zero.

When algebraic dimension is positive, we first construct an actual
canonical pullback model over a projective base. The construction
uses sufficiently large ample twists, whose sections are already
known, and steps of a single empty-boundary program. Relative
rationality allows the twist to be chosen separately at each stage
without changing that program. The base has dimension at most three.
For a surface base, fiber powers convert subadditivity into the
intersection inequality needed for nonvanishing. For a threefold
base, the genus-one Hodge line and two modular forms give an
explicit crepant klt adjoint on the base. Their common divisorial
order is determined by one integrability calculation, including
exceptional base valuations.

The algebraic-dimension-zero argument has four parts.
\begin{enumerate}
\item A nef ordinary klt adjoint has a metric of minimal
singularities with zero Lelong numbers. The proof combines a
volume-normalized capacity estimate, a differentiated
Monge--Amp\`ere equation and a Bochner transport estimate. The
comparison at the end controls a concentrating residual measure
on the same sublevel set.
\item If the divisorial locus is empty and no signed canonical
frame exists, holomorphic forms produce a transverse spherical
structure. A fixed point on the full compact convex set of positive
currents and the compactification of its holonomy lead to a
contradiction with algebraic dimension zero.
\item Relative analytic constructions produce nef reduced-boundary
models while retaining actual line identities and the global
reflexive-sheaf condition. The required dlt special termination is
proved in the dimension order used by the construction.
\item Adjunction and gluing give a section on the \emph{whole}
reduced floor. In the nonprojective torsion case, the ambient
K\"ahler class controls the pluricanonical scalars around gluing
cycles. This extra structure is what makes the cycle argument
finite.
\end{enumerate}
Finally, finite divisorial support and hard Lefschetz with multiplier
ideals extend a high power of the floor section if ambient
nonvanishing were to fail. Two choices of reduced boundary, according
to whether a signed meromorphic pluricanonical tensor exists, finish
canonical nonvanishing. The result then returns to the original nef
endpoint of the stipulated MMP.

Several of these constructions have a scope beyond their immediate
use here: the metric lemma is dimension independent; the relative
descent argument treats global reflexive sheaves directly; the
genus-one calculation determines actual crepant orders; and the
floor argument isolates the role of a single ambient K\"ahler
class in nonprojective gluing. The long projective appendices are
organized separately so that the nonprojective argument can be read
with their precise theorem statement, while every new argument used
by that theorem remains available in the same manuscript.

Figure~\ref{fig:proof-map} records these uses and the return to the
original nef endpoint. In the projective appendices,
Assumption~\ref{proj:mmp:lower-models} is the lower-dimensional induction
hypothesis, discharged in Appendix~\ref{proj:sec:completion}.
\begin{figure}[!t]
\centering
\begin{tikzpicture}[
  box/.style={draw,thin,align=center,font=\small,
    text width=4.15cm,inner sep=5pt},
  wide/.style={box,text width=10.5cm},
  flow/.style={-{Stealth[length=2mm]},thin}]
\node[wide] (endpoint) at (0,0)
  {Assumption~\ref{hyp:mmp}, retaining Assumption~\ref{hyp:campana}\\
   Original nef endpoint \((Y,\Delta)\), \(J=K_Y+\Delta\)};
\node[box,anchor=north] (projective) at (-5,-1.9)
  {Projective case\\
   Lemma~\ref{red:log-iitaka} and Theorem~\ref{proj:thm:main}\\
   Corollary~\ref{red:projective}};
\node[box,anchor=north] (direct) at (0,-1.9)
  {Uniruled or irregular case\\
   Lemmas~\ref{red:uniruled} and~\ref{red:irregular}};
\node[box,anchor=north] (canonical) at (5,-1.9)
  {Remaining smooth canonical case\\
   \(W\) non-uniruled, non-Moishezon, \(q(W)=0\)\par\smallskip
   \(a(W)>0\): Proposition~\ref{pos:nonvanishing}\\
   \(a(W)=0\): Theorem~\ref{end:canonical}, using Sections~\ref{met:section}--\ref{floor:sec}};
\node[wide] (nonvanishing) at (0,-6.3)
  {Theorem~\ref{thm:nonvanishing}: \(\kappa(Y,J)\geq0\)\\
   Nonvanishing on the original nef pair};
\node[wide] (good) at (0,-8.1)
  {Theorem~\ref{thm:main}: a Cartier multiple of \(J\)
   is globally generated on \(Y\)};
\draw[flow] (endpoint.south) -- (0,-1.05) -- (-5,-1.05) -- (projective.north);
\draw[flow] (0,-1.05) -- (direct.north);
\draw[flow] (0,-1.05) -- (5,-1.05) -- (canonical.north);
\node[fill=white,inner sep=2pt,font=\small] at (0,-1.05)
  {Section~\ref{red:section}; Proposition~\ref{red:reduction}};
\draw[flow] (projective.south) -- (-5,-5.6) -- (nonvanishing.north west);
\draw[flow] (direct.south) -- (nonvanishing.north);
\draw[flow] (canonical.south) -- (5,-5.6) -- (nonvanishing.north east);
\draw[flow] (nonvanishing.south) -- node[right,font=\small]
  {Final use of Assumption~\ref{hyp:effective-abundance}} (good.north);
\end{tikzpicture}
\caption{Reading map under Assumptions~\ref{hyp:campana}--\ref{hyp:effective-abundance}.
Lemma~\ref{red:log-iitaka} derives the sole premise of Theorem~\ref{proj:thm:main}
from Assumption~\ref{hyp:campana},
so the projective appendices add no independent conditional input.
The groupings indicate uses in this proof, not the full scope of the
component lemmas. Assumption~\ref{hyp:effective-abundance} also enters
intermediate constructions. All routes return nonvanishing to the
original nef pair before its final application on that endpoint.}
\label{fig:proof-map}
\end{figure}

